\chapter*{Appendix B: Implementation Details}
\addcontentsline{toc}{chapter}{Appendix B: Implementation Details}
\begingroup
\justifying
\setlength{\parindent}{0pt}
\setstretch{2}
\setlength{\parskip}{0.5\baselineskip}

\section{Environment Configuration}

All experiments were conducted using the following environment configuration:

\begin{verbatim}
- Python 3.10
- PyTorch 2.0+
- transformers >= 4.36.0
- trl (Transformers Reinforcement Learning)
- accelerate
- NVIDIA GPU (CUDA 11.8+)
\end{verbatim}

\section{Code Structure}

The main components of the implementation are organized as follows:

\begin{itemize}
    \item \textbf{test\_ppo\_20.py}: Main script for evaluating the PPO fine-tuned model on the test set.
    \item \textbf{train\_alignment\_discriminator.py}: Training script for the alignment discriminator.
    \item \textbf{train\_fluency\_discriminator.py}: Training script for the fluency discriminator.
    \item \textbf{ppo\_tuning\_trl.py}: PPO fine-tuning script using the TRL library.
    \item \textbf{utils/}: Utility functions for data processing, evaluation metrics, and model loading.
\end{itemize}

\section{Dataset Format}

The dataset is stored in JSON format with the following structure:

\begin{verbatim}
{
    "prompt": "Task description...",
    "error_code": "the erroneous code snippet",
    "correct_code": "the corrected code",
    "labels": ["error_type_1", "error_type_2"],
    "language": "python"
}
\end{verbatim}

\section{Hyperparameter Sensitivity}

We conducted a sensitivity analysis on key hyperparameters, including the learning rate, batch size, and PPO clipping ratio. The results indicate that the learning rate is the most sensitive hyperparameter, with values around $1 \times 10^{-5}$ yielding the best performance. The PPO clipping ratio of $\epsilon = 0.2$ was found to provide a good balance between stability and exploration.

\endgroup
